\section{Gaussian heat dissipation}\label{sec:heat}

We use Gaussian weights to turn the matrix costs of the preceding section
into derivatives of an energy.  Integration over a plane of Gaussian
centers then gives two complementary estimates: strict dissipation of the
sum of the three star energies when the variances are comparable, and
single-edge dissipation for arbitrary positive variance rates.  Both
estimates are uniform in the matrix dimension.  We retain arbitrary
positive rates in the differential identities, since both forms of
dissipation will be needed below.  Continuous Gaussian
telescoping for entangled forms also appears in
Durcik~\cite[Lemma~3 and Section~3]{Durcik2015}.

Fix $h>0$, an integer $N\ge0$, and nodes $u_i=ih$, $-N\le i\le N$.
The vertex index $v$ is taken in $\mathbb Z/3\mathbb Z$.
For finite complex edge arrays $A_v(i,j)$, put
\[
 n_v=\left(\sum_{i,j=-N}^N h^2|A_v(i,j)|^3\right)^{1/3}.
\]
For centers $p_v\in\R$ and variances $\sigma_v>0$, define
\begin{equation}\label{eq:weights}
 \begin{gathered}
 g_\sigma(t)=(2\pi\sigma)^{-1/2}e^{-t^2/(2\sigma)},
 \qquad w_v(i)=h g_{\sigma_v}(u_i-p_v),\\
 W_v(i,j)=A_v(i,j)\sqrt{w_v(i)w_{v+1}(j)}.
 \end{gathered}
\end{equation}
Form the stars and their Grams as in
\eqref{eq:row-star}--\eqref{eq:star-grams}, and set
$e_v=\Phi(T_v)=\Phi(S_v)$.
To distinguish the two blocks of the star, write
\[
 C_{v,v+1}=W_v,\qquad C_{v,v-1}=W_{v-1}^*,
 \qquad R_v=[C_{v,v+1}\ C_{v,v-1}].
\]
The diagonal Gaussian scores and the corresponding inserted Grams are
\[
 N_v=\operatorname{diag}_i\frac{u_i-p_v}{\sigma_v},\qquad
 U_v=R_v^*N_vR_v,\qquad
 P_{v,w}=C_{v,w}N_wC_{v,w}^*\quad(w\ne v).
\]
Their quadratic costs are
\begin{equation}\label{eq:heat-costs}
 V_v=b_{T_v}(U_v),\qquad
 Y_{v,w}=b_{S_v}(P_{v,w}),\qquad
 Z_v=\mathcal B_{S_v}(P_{v,v+1},P_{v,v-1}).
\end{equation}
All these insertions are supported on their respective Grams: an
inserted Gram $B^*DB$ annihilates $\ker B$ on both sides.

For $d\in\R$, equip the plane
$\Pi_d=\{p_0+p_1+p_2=d\}$ with the measure
\begin{equation}\label{eq:heat-plane-measure}
 \int_{\Pi_d}H\,\dd\pi_d
 =\int_{\R^2}H(p_0,p_1,d-p_0-p_1)\,\dd p_0\,\dd p_1.
\end{equation}
Any pair of centers can be the free coordinates, since the changes
between these charts have absolute determinant one.
Given positive rates $\boldsymbol\alpha=(\alpha_0,\alpha_1,\alpha_2)$,
we evaluate the weights at $\sigma_i=\alpha_i s$ and write
\[
 J_v(d,s)=\int_{\Pi_d}e_v\,\dd\pi_d,
 \qquad J(d,s)=\sum_vJ_v(d,s),
 \qquad \Sigma_\alpha=\sum_i\alpha_i.
\]
The arrays remain fixed whenever a scale derivative is taken.

\subsection{Differentiation and the plane identities}

\begin{lemma}[Heat identities]\label{lem:heat-identities}
For fixed finite arrays and arbitrary positive rates, the energies $e_v$
are smooth in the centers and positive variances.  For $s>0$,
\begin{align}
 2\partial_s e_v
 &=\sum_i\alpha_i\partial_{p_i}^2e_v
   -\alpha_vV_v-\sum_{w\ne v}\alpha_wY_{v,w},
   \label{eq:heat-identity}\\
 \partial_{p_{v+1}}\partial_{p_{v-1}}e_v&=Z_v.
   \label{eq:heat-mixed}
\end{align}
The plane integrals are finite, continuously differentiable in $s$, and
twice continuously differentiable in $d$, with
\begin{equation}\label{eq:heat-plane-identities}
 \begin{aligned}
 2\partial_sJ_v(d,s)
 &=\Sigma_\alpha\partial_d^2J_v(d,s)
   -\int_{\Pi_d}\left(\alpha_vV_v+
                  \sum_{w\ne v}\alpha_wY_{v,w}\right)\dd\pi_d,\\
 \partial_d^2J_v(d,s)&=\int_{\Pi_d}Z_v\,\dd\pi_d.
 \end{aligned}
\end{equation}
On compact positive scale intervals, the plane integrals of the absolute
values of all derivatives used here grow at most polynomially in $|d|$.
\end{lemma}

\begin{proof}
We first check that the possible kernels of the Grams cause no
differentiability problem.  Each star has the form
\[
 R_v=D_{\rm row}R_v^0D_{\rm col},\qquad
 R_v^0=[A_v\ A_{v-1}^*],
\]
where both diagonal multipliers are positive and invertible.  With the
column parameters fixed, $\ker R_v$ is fixed, so $T_v$ has fixed
support as the row parameters vary.  With the row parameters fixed,
$\operatorname{ran}R_v$ is fixed, so $S_v$ has fixed support as either
or both column parameters vary.  For joint smoothness, factor a nonzero
$R_v^0$ as $XY^*$ with $X,Y$ of full column rank.  Set
$X'=D_{\rm row}X$, $Y'=D_{\rm col}Y$, $P'=X'^*X'$, and $Q'=Y'^*Y'$.
These last two matrices are positive definite and smooth, and the
nonzero eigenvalues of $R_vR_v^*$ are those of
$Q'^{1/2}P'Q'^{1/2}$.  Consequently
\[
 e_v=\Phi(Q'^{1/2}P'Q'^{1/2})
\]
is smooth.  The rank-zero case is immediate.

For one Gaussian density $w=h g_\sigma(u-p)$ and its score
$N=(u-p)/\sigma$,
\[
 \partial_pw=Nw,\qquad
 \partial_p^2w=(N^2-\sigma^{-1})w,\qquad
 2\partial_\sigma w=\partial_p^2w.
\]
With the columns fixed, $T_v$ is linear in the row density, and hence
\[
 \partial_{p_v}T_v=U_v,\qquad
 2\partial_{\sigma_v}T_v=\partial_{p_v}^2T_v.
\]
The chain rule on its fixed support and Lemma~\ref{lem:hessian} give
\[
 2\partial_{\sigma_v}e_v=\partial_{p_v}^2e_v-V_v.
\]
For a column vertex $w\ne v$, use $S_v$ instead.  It is linear in that
column density, has fixed support, and satisfies
$\partial_{p_w}S_v=P_{v,w}$.  Thus
$2\partial_{\sigma_w}e_v=\partial_{p_w}^2e_v-Y_{v,w}$.
Taking the variance derivative along $\sigma_i=\alpha_i s$ proves
\eqref{eq:heat-identity}.  The two summands of
$S_v=C_{v,v+1}C_{v,v+1}^*+C_{v,v-1}C_{v,v-1}^*$ depend on different
column densities.  Its mixed derivative in those two centers is zero;
the Hessian term in the chain rule is therefore exactly $Z_v$, proving
\eqref{eq:heat-mixed}.

Before integrating these identities, we give the domination needed for
differentiation under the integral.  Lemma~\ref{lem:insertion} yields
\[
 V_v\le\tfrac34\|N_v\|_{\op}^2e_v,
 \qquad
 Y_{v,w}\le\tfrac34\|N_w\|_{\op}^2e_v,
 \qquad
 2|Z_v|\le Y_{v,v+1}+Y_{v,v-1}.
\]
For the second inequality, apply the insertion bound to $R_v^*$,
padding $N_w$ by zero on the other column block.
For any Hermitian $B$ on the row space,
\[
 \big|\tr(\sqrt{T_v}R_v^*BR_v)\big|
 \le\|B\|_{\op}e_v,
\]
because $R_v\sqrt{T_v}R_v^*$ is positive and has trace $e_v$.
The corresponding estimate for a column insertion follows by using
$S_v$.  The Gaussian derivative formulas and the Hessian chain rule
therefore bound every first center derivative, pure second center
derivative, mixed column derivative, and first scale derivative used
above by a polynomial in the centers times $e_v$, locally uniformly in
$s>0$.

The energy itself satisfies
\begin{equation}\label{eq:heat-energy-bound}
 e_v\le\left(\|W_v\|_{\HS}^2+
                  \|W_{v-1}\|_{\HS}^2\right)^{3/2}
 \le\sqrt2\left(\|W_v\|_{\HS}^3+
                       \|W_{v-1}\|_{\HS}^3\right).
\end{equation}
Since the node set is finite, each edge norm cubed is bounded by a
Gaussian in its two endpoint centers, locally uniformly for positive
variances.  On $\Pi_d$, with $|d|\le D$, each pair of distinct indices
$i,j$ satisfies
\[
 p_0^2+p_1^2+p_2^2\le4(p_i^2+p_j^2)+3D^2.
\]
Thus the preceding bounds have integrable majorants in every plane
chart, locally uniformly in $(d,s)$.  Dominated convergence justifies
all the asserted derivatives.  To obtain polynomial growth in $d$,
use the two endpoints of each bounding Gaussian as free coordinates;
the third center is $d$ minus their sum.  Integrating a polynomial times
that Gaussian leaves at most polynomial growth in $|d|$.  Constants in
these differentiability arguments may depend on the fixed arrays,
mesh, rates, and compact scale interval, but do not enter the estimates
below.

Choosing $p_i$ as the dependent coordinate gives, for each $i$,
\[
 \partial_d^2J_v=\int_{\Pi_d}\partial_{p_i}^2e_v\,\dd\pi_d.
\]
To obtain the mixed derivative, first choose $p_{v+1}$ dependent and
differentiate once.  Reparametrize the resulting integral with
$p_{v-1}$ dependent and differentiate once more.  This gives
\[
 \partial_d^2J_v
 =\int_{\Pi_d}\partial_{p_{v-1}}\partial_{p_{v+1}}e_v\,\dd\pi_d
 =\int_{\Pi_d}Z_v\,\dd\pi_d.
\]
Integrating \eqref{eq:heat-identity} proves the first identity in
\eqref{eq:heat-plane-identities}.  These comparisons concern integrated
derivatives; no pointwise equality of pure and mixed center derivatives
is asserted.
\end{proof}

\subsection{Strict dissipation at comparable variances}

The plane identities contain mixed terms $Z_v$ whose signs are not
controlled.  We bound them by column costs and then compare those costs
with the row cost at the neighboring vertex.  The resulting loss is
$\sqrt2$, small enough to leave strict dissipation for the variance
ratios we use.

\begin{lemma}[Plane dissipation]\label{lem:plane-dissipation}
For fixed finite arrays and arbitrary positive rates,
\begin{equation}\label{eq:heat-plane-dissipation}
 2\partial_sJ(d,s)\le\int_{\Pi_d}\sum_w
 \left[-\alpha_w+
 \sqrt2\max\{0,\Sigma_\alpha/2-\alpha_w\}\right]V_w\,\dd\pi_d.
\end{equation}
In particular, put $\lambda=11/10$.  For each permutation of
$(1,\lambda,\lambda)$,
\begin{equation}\label{eq:heat-strict}
 \begin{gathered}
 -\partial_sJ(d,s)\ge c_*\int_{\Pi_d}\sum_wV_w\,\dd\pi_d,\\
 c_* =\tfrac12\min\{1-\tfrac35\sqrt2,
                      \tfrac{11}{10}-\tfrac12\sqrt2\}>0.
 \end{gathered}
\end{equation}
\end{lemma}

\begin{proof}
We first prove the neighbor comparison
\begin{equation}\label{eq:heat-neighbors}
 Y_{w+1,w}+Y_{w-1,w}\le\sqrt2\,V_w.
\end{equation}
Write $A=W_w$, $B=W_{w-1}^*$, and $N=N_w$.  The elementary inequality
$\|Az+Bz'\|^2\le2\|Az\|^2+2\|Bz'\|^2$ gives
\[
 T_w\le2Q,\qquad Q=\operatorname{diag}(A^*A,B^*B).
\]
By Lemma~\ref{lem:order},
$V_w\ge2^{-1/2}b_Q([A\ B]^*N[A\ B])$.
In a block-preserving eigenbasis of $Q$, the Hessian cost is a sum of
nonnegative squared entries.  Discarding the off-diagonal blocks gives
\[
 b_Q([A\ B]^*N[A\ B])
 \ge b_{A^*A}(A^*NA)+b_{B^*B}(B^*NB).
\]
Now $A^*A\le S_{w+1}$ and $B^*B\le S_{w-1}$; the respective insertions
are $P_{w+1,w}$ and $P_{w-1,w}$.  A second use of
Lemma~\ref{lem:order}, with comparison factor one, proves
\eqref{eq:heat-neighbors}.

In \eqref{eq:heat-plane-identities}, replace
$2Z_v$ by its upper bound $Y_{v,v+1}+Y_{v,v-1}$ and sum over $v$.
The two column costs carrying the score at vertex $w$ have common
coefficient $\Sigma_\alpha/2-\alpha_w$.  If this coefficient is
negative, discard those terms; otherwise apply
\eqref{eq:heat-neighbors}.  This proves
\eqref{eq:heat-plane-dissipation}.
For the specified rates, $\Sigma_\alpha=16/5$, and the two possible
coefficients are $-1+(3/5)\sqrt2$ and $-11/10+(1/2)\sqrt2$.
Both are negative, giving \eqref{eq:heat-strict}.
\end{proof}

\subsection{Energy bounds and the single-edge case}

The decrease in \eqref{eq:heat-strict} controls the accumulated costs by
an initial energy.  We now bound that energy uniformly once the
Gaussian variances are no smaller than the mesh scale.

\begin{lemma}[Initial energy]\label{lem:initial-energy}
For all $d\in\R$ and all variances $\sigma_v\ge h^2$,
\begin{equation}\label{eq:initial-energy}
 \begin{gathered}
 \int_{\Pi_d}\|W_v\|_{\HS}^3\,\dd\pi_d\le M n_v^3,
 \qquad M=1+\sqrt{2/\pi},\\
 \int_{\Pi_d}\sum_ve_v\,\dd\pi_d
       \le2\sqrt2 M\sum_v n_v^3.
 \end{gathered}
\end{equation}
These bounds also hold with the original right-hand sides if any entries
are replaced by entries of smaller absolute value.
\end{lemma}

\begin{proof}
The Gaussian grid mass obeys
\begin{equation}\label{eq:heat-mesh-mass}
 \sum_{i=-N}^Nh g_\sigma(u_i-p)
 \le1+\frac{2h}{\sqrt{2\pi\sigma}}.
\end{equation}
Indeed, for $u_i\ge p+h$ compare the summand with the integral over the
preceding interval of length $h$; for $u_i\le p-h$, use the succeeding
interval.  These intervals are disjoint, and the at most two remaining
nodes contribute at most $2h\sup g_\sigma$.
For $\sigma\ge h^2$, the right side is at most $M$.
Weighted H\"older therefore gives
\[
 \|W_v\|_{\HS}^3
 =\left(\sum_{i,j}|A_v(i,j)|^2w_v(i)w_{v+1}(j)\right)^{3/2}
 \le M\sum_{i,j}|A_v(i,j)|^3w_v(i)w_{v+1}(j).
\]
Use $p_v,p_{v+1}$ as free coordinates on $\Pi_d$.  Each density
integrates to $h$, proving the first estimate in
\eqref{eq:initial-energy}.  Each edge occurs in two stars, so
\eqref{eq:heat-energy-bound} proves the second.  The same upper bounds
decrease when the entrywise absolute values decrease.
\end{proof}

For example, if $h^2\le s_0<s_1$ and the rates are a permutation of
$(1,\lambda,\lambda)$, these two lemmas imply
\begin{equation}\label{eq:heat-fixed-array-integral}
 \int_{s_0}^{s_1}\int_{\Pi_d}\sum_vV_v\,\dd\pi_d\,\dd s
 \le\frac{J(d,s_0)-J(d,s_1)}{c_*}
 \le\frac{2\sqrt2 M}{c_*}\sum_v n_v^3.
\end{equation}
The first inequality retains the terminal energy, which is essential
when arrays are changed at finitely many scales.

For much more unequal variance rates,
\eqref{eq:heat-plane-dissipation} does not guarantee that the summed
star energy decreases.  A single edge nevertheless has no mixed term
and retains exact dissipation.  We state that consequence separately.

\begin{lemma}[Single-edge energy]\label{lem:single-edge-energy}
Fix $v\ne w$, retain the weighted block $C=C_{v,w}$, and delete the
other block of the star at $v$.  For arbitrary positive rates, set
\[
 \begin{aligned}
 e^e_{v,w}&=\Phi(CC^*),\qquad
 V^e_{v,w}=b_{C^*C}(C^*N_vC),\qquad
 Y^e_{v,w}=b_{CC^*}(CN_wC^*),\\
 J^e_{v,w}(s)&=\iint_{\R^2}e^e_{v,w}\,\dd p_v\,\dd p_w.
 \end{aligned}
\]
The plane integrals of $e^e_{v,w}$, $V^e_{v,w}$, and $Y^e_{v,w}$ are
independent of $d$.
Moreover,
\begin{equation}\label{eq:heat-single-edge}
 -2\partial_sJ^e_{v,w}(s)
 =\iint_{\R^2}
      (\alpha_vV^e_{v,w}+\alpha_wY^e_{v,w})\,\dd p_v\,\dd p_w.
\end{equation}
If $n_e$ denotes the discrete $L^3$ norm of this edge's unweighted
array, then, whenever $\alpha_vs_0,\alpha_ws_0\ge h^2$ and $s_1>s_0$,
\begin{equation}\label{eq:heat-single-edge-integral}
 \int_{s_0}^{s_1}\iint_{\R^2}
      (\alpha_vV^e_{v,w}+\alpha_wY^e_{v,w})
          \,\dd p_v\,\dd p_w\,\dd s
 \le2M n_e^3.
\end{equation}
For the original two-block star, $Y_{v,w}\le Y^e_{v,w}$.
\end{lemma}

\begin{proof}
The three integrands depend only on the two endpoint centers.
Using these as free coordinates proves independence of $d$.
In the one-edge star the other column insertion is zero, so $Z_v=0$.
After permuting column blocks, its matrices are
\[
 R=[C\ 0],\qquad T=\operatorname{diag}(C^*C,0),\qquad
 U=\operatorname{diag}(C^*N_vC,0),\qquad S=CC^*.
\]
The Hessian cost is unchanged by zero extension, so its row and column
costs are exactly $V^e_{v,w}$ and $Y^e_{v,w}$.
Equation~\eqref{eq:heat-plane-identities} now proves
\eqref{eq:heat-single-edge}.
Also $e^e_{v,w}\le\|C\|_{\HS}^3$, so
Lemma~\ref{lem:initial-energy} gives $J^e_{v,w}(s_0)\le M n_e^3$.
Integrate \eqref{eq:heat-single-edge} and discard the nonnegative
terminal energy to obtain \eqref{eq:heat-single-edge-integral}.
Finally $CC^*\le S_v$, while $CN_wC^*$ is supported on $CC^*$.
Lemma~\ref{lem:order} with factor one gives $Y_{v,w}\le Y^e_{v,w}$.
\end{proof}
